\section{TEM Specimen Tilt Navigation}

This note fixes the conventions, the geometry and the identifiability analysis behind
\texttt{pytex.tem}: given a crystal at a known holder position with an indexed zone axis, determine
the holder tilts $(\alpha, \beta)$ that bring a requested zone axis onto the electron beam.

The forward problem is a single rigid-body equation and the inverse problem has a two-line closed
form. Almost all of the difficulty lies elsewhere, in what the available observations do
\emph{not} determine, and the bulk of this note is therefore devoted to that question.

\subsection{Frames}

No new frame domain is introduced. Per the repository frame doctrine the holder frame \emph{is}
the specimen-domain frame for TEM work, so the crystal-to-holder rotation is an ordinary
crystal-to-specimen \texttt{Orientation}.

\begin{itemize}
  \item \textbf{Crystal} $C$: right-handed Cartesian tied to the lattice through the structure
    matrix $\mathbf{A}$, with $[uvw] \mapsto \mathbf{A}[u,v,w]^{\mathsf T}$ and
    $(hkl) \mapsto \mathbf{A}^{*}[h,k,l]^{\mathsf T}$, $\mathbf{A}^{*} = (\mathbf{A}^{-1})^{\mathsf T}$.
  \item \textbf{Holder} $H$: rigidly attached to the specimen and to the cartridge, with
    $\hat{\mathbf{x}}_H$ along the rod (the $\alpha$ axis), $\hat{\mathbf{y}}_H$ along the nominal
    $\beta$ axis. Coincides with the laboratory frame at zero tilt.
  \item \textbf{Laboratory} $L$: fixed to the column, $\hat{\mathbf{z}}_L$ pointing up the column
    toward the gun, so electrons propagate along $-\hat{\mathbf{z}}_L$.
  \item \textbf{Pattern} $P$: axes of the stored diffraction-image array, origin at the
    transmitted beam.
\end{itemize}

\subsection{The Specimen/Holder Collapse}

A textbook treatment inserts a specimen frame between crystal and holder, joined by a mounting
rotation $\mathbf{M}_{H \leftarrow S}$. \textbf{That rotation and the crystal orientation are not
separately identifiable from diffraction data.} No observable in this problem depends on either
alone, only on the product, so carrying them separately manufactures an unobservable degree of
freedom. PyTex therefore carries the single crystal-to-holder rotation

\begin{equation}
\mathbf{U} \;\equiv\; \mathbf{M}_{H \leftarrow S}\,\mathbf{M}_{S \leftarrow C}.
\label{eq:tn-collapse}
\end{equation}

The mounting rotation reappears only in the tilt envelope, where a grid bar shadows at a
mounting-dependent azimuth, never as an orientation unknown.

\subsection{Stage Kinematics}

The two tilt axes of a double-tilt holder are not both fixed. The $\alpha$ axis is the rod, fixed
in the column; the $\beta$ axis is a cradle carried \emph{inside} the rod and therefore moves with
$\alpha$. Writing the $\beta$ rotation about its instantaneous laboratory-frame axis
$\mathbf{R}_x(\alpha)\hat{\mathbf{y}}$ and composing with $\alpha$,

\begin{equation}
\mathbf{R}_{\text{stage}}
= \Bigl[\mathbf{R}_x(\alpha)\,\mathbf{R}_y(\beta)\,\mathbf{R}_x(\alpha)^{\mathsf T}\Bigr]\,
  \mathbf{R}_x(\alpha)
= \mathbf{R}_x(\alpha)\,\mathbf{R}_y(\beta),
\label{eq:tn-stage}
\end{equation}

so the moving-axis factors cancel exactly. The cancellation is a property of this particular axis
pair and \emph{not} a general licence to ignore axis motion: it does not survive a non-orthogonal
or mis-set pair, which is why the calibrated model composes the two rotations explicitly.

The master equation is then

\begin{equation}
\mathbf{v}_L(\alpha,\beta) = \mathbf{R}_{\text{stage}}(\alpha,\beta)\,\mathbf{U}\,\mathbf{v}_C .
\label{eq:tn-master}
\end{equation}

Everything below is an inversion of, or a constraint on, \eqref{eq:tn-master}.

\subsection{Beam Direction, And Why It Is The Useful Form}

Transposing \eqref{eq:tn-master} for the beam axis gives

\begin{equation}
\hat{\mathbf{b}}_H(\alpha,\beta)
= \mathbf{R}_{\text{stage}}^{\mathsf T}\hat{\mathbf{z}}_L
= \bigl(-\cos\alpha\,\sin\beta,\ \sin\alpha,\ \cos\alpha\,\cos\beta\bigr).
\label{eq:tn-beam}
\end{equation}

This is a spherical coordinate system whose \textbf{pole is the $\beta$ axis}, and three
consequences follow immediately.

\begin{enumerate}
  \item The Jacobian is $\cos\alpha$, so a tilt rectangle reaches the solid angle
    \begin{equation}
    \Omega = (\beta_{\max} - \beta_{\min})\,(\sin\alpha_{\max} - \sin\alpha_{\min}).
    \label{eq:tn-solid-angle}
    \end{equation}
    A $\pm 30^{\circ}$ holder gives $\Omega = \pi/3 = 1.0472\,\mathrm{sr}$, which is
    $8.33\,\%$ of all directions.
  \item Curves of constant $\beta$ are great circles through $\pm\hat{\mathbf{y}}_H$, and curves of
    constant $\alpha$ are small circles about it. The accessible region is a geodesic quadrilateral
    and draws exactly, as circular arcs.
  \item $\lVert \partial\hat{\mathbf{b}}_H/\partial\alpha \rVert = 1$ but
    $\lVert \partial\hat{\mathbf{b}}_H/\partial\beta \rVert = \cos\alpha$: a degree of $\beta$ buys
    only $\cos\alpha$ degrees of crystal rotation. This is the double-tilt holder's gimbal lock.
\end{enumerate}

\subsection{The Closed-Form Solution}

Let $\hat{\mathbf{w}} = \mathbf{U}\hat{\mathbf{t}}_C = (w_1, w_2, w_3)^{\mathsf T}$ be the target
direction in holder coordinates. Alignment requires $\hat{\mathbf{w}} = \pm\hat{\mathbf{b}}_H$,
and matching against \eqref{eq:tn-beam} gives

\begin{equation}
\beta^{*} = \operatorname{atan2}(-w_1,\ w_3),
\qquad
\alpha^{*} = \operatorname{atan2}\!\left(w_2,\ \sqrt{w_1^2 + w_3^2}\right).
\label{eq:tn-solution}
\end{equation}

Verification: with $\rho = \sqrt{w_1^2 + w_3^2}$, $\mathbf{R}_y(\beta^{*})^{\mathsf T}$ carries
$\hat{\mathbf{w}}$ to $(0, w_2, \rho)$ and $\mathbf{R}_x(\alpha^{*})^{\mathsf T}$ carries that to
$(0,0,1)$.

Four branches satisfy the alignment condition. Two require $\lvert\beta\rvert > 90^{\circ}$ and one
requires roughly $180^{\circ}$ of $\alpha$, so on a real holder exactly one is relevant.
\textbf{The combinatorial richness of the answer therefore comes from symmetry and ambiguity, not
from the trigonometry.} When $\rho = 0$ the target lies along the $\beta$ axis, $\beta$ is
indeterminate and $\alpha = \pm 90^{\circ}$; the degeneracy is detected rather than producing an
arbitrary $\operatorname{atan2}(0,0)$.

For a calibrated stage with a non-orthogonal or mis-set axis pair there is no closed form, and
\eqref{eq:tn-solution} seeds a Gauss--Newton refinement on the two-dimensional residual.

\subsection{Reconstruction, And Why Two Zones Beat One Pattern}

\subsubsection{Single pattern}

Pattern indexing supplies the crystal-to-pattern rotation $\mathbf{R}_{P \leftarrow C}$, and

\begin{equation}
\mathbf{U} = \mathbf{R}_{\text{stage}}(\alpha_c,\beta_c)^{\mathsf T}\,
             \mathbf{F}\,\mathbf{R}_z(\varphi_D)\,\mathbf{R}_{P \leftarrow C},
\label{eq:tn-mode-a}
\end{equation}

with $\varphi_D$ the diffraction rotation and $\mathbf{F}$ the parity of the stored image. If
$\mathbf{F}$ is wrong the product is improper, so $\det \mathbf{U} = -1$ is a free and exact
self-check.

\subsubsection{Two zone axes}

If zone axes $\hat{\mathbf{n}}_1, \hat{\mathbf{n}}_2$ were indexed at two stage positions then by
\eqref{eq:tn-beam}

\begin{equation}
\mathbf{U}\,\hat{\mathbf{n}}_i = \hat{\mathbf{b}}_H(\alpha_i, \beta_i), \qquad i = 1,2,
\label{eq:tn-mode-b}
\end{equation}

and two non-parallel correspondences determine a rotation uniquely (Wahba's problem). \textbf{This
path requires neither $\varphi_D$ nor $\mathbf{F}$}: the single hardest calibration constant in the
problem is not needed if the operator has visited two zones, which after any real session of
Kikuchi-band chasing they almost always have.

It also yields a calibration-free consistency test. The interzonal angle is a crystallographic
invariant while the beam-direction angle depends only on the stage model, so

\begin{equation}
\Delta_{\text{cons}} = \Bigl\lvert\, \angle(\hat{\mathbf{n}}_1,\hat{\mathbf{n}}_2)
- \angle\bigl(\hat{\mathbf{b}}_H(\alpha_1,\beta_1), \hat{\mathbf{b}}_H(\alpha_2,\beta_2)\bigr) \Bigr\rvert
\label{eq:tn-consistency}
\end{equation}

indicts, in order of likelihood, a wrong zone-axis sign, a wrong stage sign convention, a wrong
indexing, or a bent specimen.

A residual two-fold ambiguity survives: flipping both zone-axis senses admits
$\mathbf{U}' = \mathbf{U}\mathbf{Q}$ with $\mathbf{Q}$ the $180^{\circ}$ rotation about
$\hat{\mathbf{n}}_1 \times \hat{\mathbf{n}}_2$. It is harmless exactly when $\mathbf{Q}$ is an
element of the crystal's proper point group, which is a one-line check. For cubic zones $[001]$ and
$[110]$ the axis is $\langle 1\bar{1}0 \rangle$ and the answer is yes.

\subsection{Identifiability: Three Independent Ambiguities}

The classical ``$180^{\circ}$ ambiguity'' is three unrelated things. Conflating them is why the
failure mode has a reputation for being intractable.

\subsubsection{Layer 1: Friedel/Laue}

Kinematic intensities obey $\lvert F(\mathbf{g})\rvert = \lvert F(-\mathbf{g})\rvert$, so the
recorded pattern is centrosymmetric even when the crystal is not. The reconstruction is determined
only up to the observation stabilizer

\begin{equation}
G_{\text{obs}} = \bigl\{\, \mathbf{Q} \in L \cap SO(3)\ :\ \mathbf{Q}\ \text{maps the zone plane to itself} \,\bigr\},
\label{eq:tn-stabilizer}
\end{equation}

with $L$ the \emph{Laue class}. The genuinely distinct hypotheses are counted by the cosets
$G_{\text{obs}} / (G_{\text{obs}} \cap G_{\text{proper}})$.

The decisive observation is that $L \cap SO(3)$ exceeds the crystal's proper point group for
\textbf{exactly ten of the thirty-two point groups} --- those containing improper operations other
than inversion --- and by a factor of exactly two in every case:

\begin{center}
\begin{tabular}{lll}
\hline
Class & Point groups & Index \\
\hline
Centrosymmetric (11) & $\bar{1}$, $2/m$, $mmm$, $4/m$, $4/mmm$, $\bar{3}$, $\bar{3}m$, $6/m$, $6/mmm$, $m\bar{3}$, $m\bar{3}m$ & 1 \\
Enantiomorphic (11) & $1$, $2$, $222$, $4$, $422$, $3$, $32$, $6$, $622$, $23$, $432$ & 1 \\
Improper, acentric (10) & $m$, $mm2$, $\bar{4}$, $4mm$, $\bar{4}2m$, $3m$, $\bar{6}$, $6mm$, $\bar{6}m2$, $\bar{4}3m$ & 2 \\
\hline
\end{tabular}
\end{center}

Two consequences. \textbf{For a centrosymmetric crystal Friedel's law adds nothing}: every
``ambiguous'' reconstruction is a symmetry-equivalent description of the same physical orientation,
so the much-feared $180^{\circ}$ ambiguity is harmless for every cubic and hexagonal metal. And
\emph{non-centrosymmetric is the wrong test}: quartz, point group $32$, is acentric yet entirely
ambiguity-free, because its point group is already all rotations.

Where the ambiguity is real, the destination is a Friedel partner: the SAED pattern there is
indistinguishable, but the physical orientation differs, and with it every polarity-sensitive
measurement --- CBED/HOLZ symmetry, polar-axis determination, $\mathbf{g}\cdot\mathbf{b}$ analysis.

\subsubsection{Layer 2: Instrumental rotation}

An error $\delta\varphi$ in the diffraction rotation is a rotation about the beam axis. It is
\emph{not} absorbed by crystal symmetry, and following the tilts it implies lands the specimen

\begin{equation}
\Delta_{\text{res}} = 2\arcsin\!\bigl(\sin(\delta\varphi/2)\,\sin\theta_c\bigr)
\ \approx\ \delta\varphi\,\sin\theta_c
\label{eq:tn-residual-law}
\end{equation}

away from the target, where $\theta_c$ is the angle between the current and target zone axes.
Derivation: with $\mathbf{c} = \mathbf{U}_{\text{true}}\hat{\mathbf{t}}_C$, the tilts satisfy
$\mathbf{R}_{\text{stage}}\mathbf{R}_z(\delta\varphi)\mathbf{c} = \hat{\mathbf{z}}_L$ while the
specimen delivers $\mathbf{R}_{\text{stage}}\mathbf{c}$; the residual is
$\angle(\mathbf{c}, \mathbf{R}_z(\delta\varphi)\mathbf{c})$.

Equation \eqref{eq:tn-residual-law} carries the two operationally important facts. At
$\delta\varphi = 180^{\circ}$ both tilt angles are negated: the specimen goes exactly the wrong way
while the calculation reports a clean zero residual. And the error \emph{scales with the length of
the hop}, so a long excursion should be routed through intermediate low-index zones with
re-indexing at each, which converts an open-loop calculation into a self-correcting procedure.

\subsubsection{Layer 3: Parity}

The image mirroring $\mathbf{F}$ and the readout signs $s_\alpha, s_\beta$ reflect the trajectory
rather than rotating it, so unlike layer 2 they cannot be absorbed into any single angle. They are
caught by $\det\mathbf{U}$ and by \eqref{eq:tn-consistency}.

\subsection{Calibration}

The instrument reports the two stage angles and nothing about the transformation between them and
the diffraction pattern. The core measurement is a two-excursion procedure: apply a known positive
$\alpha$, record the azimuth $\psi_\alpha$ at which a tracked Kikuchi feature moved, repeat for
$\beta$.

A rigid rotation carries a pole at the pattern centre to $\mathbf{R}\hat{\mathbf{z}}_L$. At zero
tilt $\mathbf{R}_x(\Delta\alpha)\hat{\mathbf{z}}_L$ displaces along $-\hat{\mathbf{y}}_L$ and
$\mathbf{R}_y(\Delta\beta)\hat{\mathbf{z}}_L$ along $+\hat{\mathbf{x}}_L$, so with
$\mathbf{R}_{L \leftarrow P} = \mathbf{F}\mathbf{R}_z(\varphi_D)$,

\begin{equation}
\varphi_D = -\psi_\beta,
\qquad
\psi_\alpha - \psi_\beta =
\begin{cases}
-90^{\circ} & \mathbf{F} = \mathbf{I},\\
+90^{\circ} & \text{mirrored}.
\end{cases}
\label{eq:tn-calibration}
\end{equation}

Three results from two exposures --- the rotation numerically, the parity from a sign, and a
redundancy check, since the azimuths must be $90^{\circ}$ apart. The displacement magnitude is a
bonus: a pole must move by exactly the applied angle, which also confirms the camera length.

$\varphi_D$ is hysteretic in the projector and diffraction lens settings, so a calibration records
the camera length and voltage it was measured at and must not be interpolated to another. A
plausible interpolated value for a hysteretic quantity is precisely the layer-2 failure above.

\subsection{Eucentric Height Is Not An Orientation Parameter}

Being off eucentric height introduces \emph{no orientation error to first order}: tilting about a
displaced axis is a rotation composed with a translation, and only the rotation enters
\eqref{eq:tn-master}. The real cost is that the region of interest translates out of the
selected-area aperture, so the operator loses the crystal rather than mis-orienting it. This
belongs in path segmentation, not in the orientation model.

\subsection{Path Planning}

A straight line in $(\alpha,\beta)$ is not a straight line on the crystal sphere. The geodesic
between two zone axes lies in the plane they span, whose normal rationalizes to a low-index
$(hkl)$ --- and the Kikuchi band of that reflection is exactly the band joining the two poles.
\textbf{The mathematically optimal path is the one experienced operators already follow by eye},
which is why the planner reports the band to follow.

\subsection{Failure Modes}

In order of practical likelihood:

\begin{enumerate}
  \item $\varphi_D$ wrong, worst at $180^{\circ}$: a clean, confident, exactly-backwards answer,
    detected only by calibration or by the two-zone path.
  \item Parity or sign convention wrong: reflected trajectory, detected by $\det\mathbf{U}$ and by
    \eqref{eq:tn-consistency}.
  \item Specimen bent or the area polycrystalline: the rigid assumption fails and no reconstruction
    is valid; bounded by \eqref{eq:tn-consistency} across pairs.
  \item Target orbit outside the envelope: not a failure but an answer, and the most common useful
    output.
  \item Target near the $\beta$-axis pole: ill-conditioned by $1/\cos\alpha$, bounded at $1.31$
    within a real envelope.
  \item One of the ten affected point groups with a polarity-sensitive purpose: irreducible from
    kinematic SAED; reported, not guessed.
\end{enumerate}

\subsection{References}

\begin{itemize}
  \item De Graef, M., \emph{Introduction to Conventional Transmission Electron Microscopy},
    Cambridge University Press, DOI: 10.1017/CBO9780511615092 --- stage geometry, diffraction
    rotation, Kikuchi navigation.
  \item Williams, D. B. and Carter, C. B., \emph{Transmission Electron Microscopy}, Springer,
    DOI: 10.1007/978-0-387-76501-3 --- indexing, tilt practice, calibration procedures.
  \item Wahba, G., \emph{SIAM Review} \textbf{7} (1965) 409, DOI: 10.1137/1007077 --- the
    two-vector attitude problem underlying the two-zone reconstruction.
  \item Hahn, Th. (ed.), \emph{International Tables for Crystallography, Volume A}, IUCr/Springer,
    DOI: 10.1107/97809553602060000100 --- point groups, Laue classes, stabilizer orders.
  \item Morawiec, A., \emph{Orientations and Rotations}, Springer, DOI: 10.1007/978-3-662-09156-2
    --- rotation representations, symmetry orbits, tangent-space covariance.
\end{itemize}
